\subsection{Mario Zambrano · PUB/SUB aplicado al SGRA}

\paragraph{Arquitectura y problema identificado.}
El SGRA permite gestionar solicitudes de refuerzo académico
mediante Angular, Spring Boot y PostgreSQL. Actualmente, el
backend registra las solicitudes y gestiona notificaciones
por correo. Como propuesta de mejora, se plantea distribuir
eventos a componentes independientes, como notificaciones,
auditoría e indicadores \cite{sgra_backend}.

\paragraph{Fundamentos de Publish/Subscribe.}
PUB/SUB utiliza publicadores, un intermediario de mensajería
y suscriptores. El publicador no necesita conocer directamente
a cada receptor. Esto facilita el desacoplamiento y permite
procesar actividades complementarias de manera asíncrona
\cite{lazidis2022}.

\paragraph{Demostración y aplicación al SGRA.}
Primero se demuestra PUB/SUB con Mosquitto: un publicador
envía el evento \texttt{SolicitudTutoriaCreada} y dos
suscriptores reciben el mismo mensaje.

Posteriormente, se presenta la integración experimental
realizada en Java. Se incorporó Eclipse Paho al backend y
se ejecutaron pruebas de publicación, recepción, caída del
broker, recuperación y eventos duplicados.

Además, tres pruebas automatizadas con Mockito verificaron
el comportamiento del servicio ante registros exitosos,
fallos del registro y errores MQTT. Una prueba adicional
comprobó la publicación desde el servicio del SGRA hacia
Mosquitto y la recepción por dos suscriptores.

\paragraph{Resultados y limitaciones.}
Las pruebas confirmaron la distribución de eventos y
permitieron observar que una caída del broker impide nuevas
publicaciones mientras permanezca desconectado.

También se comprobó que dos publicaciones con el mismo
identificador lógico llegan a los suscriptores. Por ello,
una implementación completa necesitaría controlar duplicados
y recuperar publicaciones fallidas.

La persistencia se simuló mediante Mockito: no se ejecutó
una transacción real en PostgreSQL. Como mejora futura se
propone evaluar Transactional Outbox y consumidores
idempotentes para reforzar la consistencia y confiabilidad
\cite{microsoft2026}.
